\documentclass{article}
\usepackage{ctex}%中文
\usepackage{amsmath}
\begin{document}
\section{bezier曲线}
为了避免高阶Bezier曲线存在的数值稳定性较差的问题对路径优化造成的影响，同时确保路径的平滑处理，本文采用低阶的Bezier曲线进行轨迹的平滑处理。在处理路径拐点时，观察到单个拐点的出现频率较高，且可能连续出现多个拐点。尽管二阶Bezier曲线可以较好地处理单个拐点的路径，但对于连续出现多个拐点的情况，其平滑处理较为困难。相比之下，三阶Bezier曲线则能更有效地应对这种情况，不会出现处理连续多拐点的困难。[三阶]

\section{碰撞检测}
外接圆碰撞检测法

包围盒碰撞检测法

计算几何法




\section{A*算法，对比}
\subsection{介绍}
算法中有一个启发式函数$h(n)$，它估计从节点$n$到目标节点的代价。A*算法通过结合实际代价$g(n)$和启发式代价$h(n)$来选择下一个要扩展的节点。A*算法的核心思想是优先扩展具有最低$f(n) = g(n) + h(n)$值的节点，其中$g(n)$是从起点到节点$n$的实际代价，$h(n)$是从节点$n$到目标节点的启发式估计代价。A*算法在路径搜索中具有较高的效率和准确性，广泛应用于各种领域，如机器人导航、游戏开发和网络路由等。
Dijkstra算法是一种经典的图算法，用于解决单源最短路径问题。它通过逐步扩展节点来找到从起点到所有其他节点的最短路径。Dijkstra算法的核心思想是使用一个优先队列来存储待扩展的节点，并根据当前已知的最短路径代价进行排序。每次从优先队列中选择具有最低代价的节点进行扩展，并更新其邻居节点的代价。Dijkstra算法保证了在没有负权边的情况下能够找到最短路径，但在存在负权边时可能会出现问题。
贪心最佳优先算法是一种启发式搜索算法，它在每一步选择当前看起来最优的选项，以期望最终能够找到全局最优解。与A*算法不同，贪心最佳优先算法只考虑启发式代价$h(n)$，而不考虑实际代价$g(n)$。因此，贪心最佳优先算法可能会陷入局部最优解，而无法保证找到全局最优解。尽管如此，贪心最佳优先算法在某些情况下仍然可以提供较快的搜索速度，特别是在启发式函数非常准确的情况下。

Dijkstra 算法和 A*算法都是基于图中已知障碍物的信息，通过对图中可通行区域来搜寻从起点到目标点的路径。其中，Dijkstra 算法采用的是贪心策略，向地图中每个方向进行探索，并保存所有搜索到的点，直到遍历到目标点，然后在探索到的点中构建出起点与目标点的路径。[停车场] A*算法则在Dijkstra算法的基础上引入了启发式函数，能够更有效地引导搜索过程，减少不必要的节点扩展，从而提高搜索效率。
\subsection{对比}
其中，A*算法在路径搜索中具有较高的效率和准确性，广泛应用于各种领域，如机器人导航、游戏开发和网络路由等。Dijkstra算法是一种经典的图算法，用于解决单源最短路径问题。它通过逐步扩展节点来找到从起点到所有其他节点的最短路径。Dijkstra算法的核心思想是使用一个优先队列来存储待扩展的节点，并根据当前已知的最短路径代价进行排序。每次从优先队列中选择具有最低代价的节点进行扩展，并更新其邻居节点的代价。Dijkstra算法保证了在没有负权边的情况下能够找到最短路径，但在存在负权边时可能会出现问题。贪心最佳优先算法是一种启发式搜索算法，它在每一步选择当前看起来最优的选项，以期望最终能够找到全局最优解。与A*算法不同，贪心最佳优先算法只考虑启发式代价$h(n)$，而不考虑实际代价$g(n)$。因此，贪心最佳优先算法可能会陷入局部最优解，而无法保证找到全局最优解。尽管如此，贪心最佳优先算法在某些情况下仍然可以提供较快的搜索速度，特别是在启发式函数非常准确的情况下。
\subsection{启发式函数}
可用的h（x）函数有很多种，常见的包括曼哈顿距离、欧几里得距离和切比雪夫距离等。选择合适的启发式函数对于A*算法的性能至关重要，因为它直接影响到搜索过程中的节点扩展顺序和效率。一个好的启发式函数应该能够准确地估计从当前节点到目标节点的代价，同时保持计算效率，以确保A*算法能够快速找到最优路径。

曼哈顿距离
适合：只能上下左右四向移动的网格

欧几里得距离
适合：可以任意方向移动（连续空间）

切比雪夫距离
适合：八向移动（上下左右 + 斜向）
\[
h_{\text{曼哈顿}}(x) = |x_1 - x_2| + |y_1 - y_2|
\]
\[
h_{\text{欧几里得}}(x) = \sqrt{(x_1 - x_2)^2 + (y_1 - y_2)^2}
\]
\[
h_{\text{切比雪夫}}(x) = \max\bigl(|x_1 - x_2|,\ |y_1 - y_2|\bigr)
\]
\subsection{环境复杂度量化}
看[1]








\end{document}